\section{What the analyzer does, end to end}

The analyzer turns four CSV files into two things: one row of numbers per run, and a fitted
polynomial per membrane with the best setup that polynomial allows. Nothing in between is
saved to disk. Every stage takes objects in and hands objects on, and the whole chain is
three statements long.

\subsection{The four stages, in the order they run}

\texttt{Analyzer.analyze} in \texttt{src/limodel/analyzer/analyzer.py} is the order, and it is
short on purpose --- the order is the only thing this class knows:

\begin{enumerate}[nosep]
  \item \textbf{Load.} \texttt{DataLoader.load} reads the four CSV files and returns a list of
        \texttt{Run} objects, each with its readings already inside it.
  \item \textbf{Preprocess.} \texttt{Preprocessor.run} takes that list and gives back the
        shorter list the modelling may use. On the way it clears single values it cannot
        believe.
  \item \textbf{Calculate.} \texttt{Calculator.find\_results} takes the surviving runs
        and returns one \texttt{Result} per run whose two outcomes came out.
  \item \textbf{Model.} \texttt{Surrogate.find\_models} takes those results, groups them by
        membrane, fits three outcomes per membrane, and searches each membrane for its best
        setup.
\end{enumerate}

Each stage sees only what the one before it returned. The preprocessor never sees a CSV file;
the calculator never sees a run that was dropped; the fitting never sees a \texttt{Run} at all,
only \texttt{Result} rows.

What comes back is an \texttt{Analysis}: the runs used, the results, and the models. The runs
travel beside the results rather than being replaced by them, because a reader checking a
number wants the readings it came from, and because the count of runs used is part of the
claim --- a polynomial from sixteen runs and one from twenty-one are different statements.

\subsection{The journal: made once, handed down}

The \texttt{Journal} is created empty in one place --- the \texttt{analyze} command in
\texttt{cli/main.py} --- and passed into \texttt{Analyze}, which passes it into
\texttt{Analyzer}, which passes it into \texttt{Preprocessor}, \texttt{Calculator} and
\texttt{Surrogate}. There is exactly one of them per run of the program.

It is injected rather than owned by each stage for one reason: the report has to add up. If the
preprocessor kept its own list of drops and the calculator kept another, the two would have to
be collected and merged by whoever printed them, and a stage added later would be easy to leave
out of that merge. With one journal, a stage that writes nothing down has simply written nothing
down, and the arithmetic at the top of the report is a count of one list.

Two things get written. A \textbf{drop} means the whole record is out of the modelling. A
\textbf{skip} means the record stays and one of its values does not. An \texttt{Entry} carries
the action, the record's id, the record's class name, the reason in plain words, and (for a
skip) the field that was cleared.

The journal takes the id and the class name off the record itself, inside
\texttt{Journal.\_write}, instead of having each caller pass them. That is deliberate: a caller
that passed its own id could pass one that disagreed with the record, and nothing would catch
it.

Because every removal leaves an entry, the line under the journal reconciles by construction:

\begin{equation}
  \label{eq:reconcile}
  N_{\text{loaded}} \;=\; N_{\text{used}} \;+\; \bigl|\{\, \text{drop ids}\,\} \setminus
  \{\, \text{used ids}\,\}\bigr|
\end{equation}

The set difference in \eqref{eq:reconcile} matters. A run can be dropped by the calculator
\emph{after} the preprocessor already passed it, so it would appear both in the used list
handed over and in the drops; subtracting the used ids counts each run once. On the current
data the analyzer writes 32 entries --- 5 drops and 27 skips --- and the report prints
\texttt{using 16 of 21 runs}.

The \texttt{Surrogate} stage writes nothing to the journal. A membrane it will not fit is
reported on its own answer, in \texttt{SurrogateModels.not\_fitted}, because that is a finding
about the design of the experiments rather than a record that was set aside.

\subsection{The four files, and why loading is all or nothing}

\begin{table}[htbp]
  \centering
  \caption{The four CSV files in \texttt{data/}, with the row counts in the current data.}
  \label{tab:files}
  \begin{tabular}{@{}llr@{}}
    \toprule
    File & One row per & Rows \\
    \midrule
    \texttt{runs.csv}           & experiment, with how it was set up        & 21 \\
    \texttt{tank\_readings.csv} & run, minute and tank                      & 1049 \\
    \texttt{stack\_readings.csv}& run and minute, from the power log        & 891 \\
    \texttt{samples.csv}        & bottle sent to the lab                    & 345 \\
    \bottomrule
  \end{tabular}
\end{table}

\texttt{runs.csv} holds the setup and nothing measured during the run: the membranes, the feed
concentrations, the set current, the flow, the membrane area, the cell pairs, the three tank
tares, and the \texttt{include} flag. \texttt{tank\_readings.csv} holds conductivity, pH,
temperature and weight for the concentrate, dilute and rinse.
\texttt{stack\_readings.csv} holds the two voltage columns and the two current columns ---
supply and multimeter --- which is what lets the checks compare one instrument against the
other. \texttt{samples.csv} holds the lab's answers: lithium, magnesium, sodium, potassium,
calcium and TOC.

Each file has a schema in \texttt{src/limodel/schema/}, written as a \emph{pandera}
\texttt{DataFrameModel}. Every schema sets \texttt{coerce = True} and \texttt{strict = True}, so
values are converted to the declared type on load and the columns must be exactly the declared
ones --- no extras, none missing. The three reading files declare a uniqueness key
(\texttt{run\_code} with \texttt{minute}, plus \texttt{tank} where there is one), and
\texttt{run\_code} and \texttt{sample\_id} are unique in their own files. Measured columns are
\texttt{nullable}; setup columns are not, and most carry a range (\texttt{li\_feed\_mg\_l}
greater than zero, \texttt{ph} between 0 and 14, \texttt{area\_cm2} greater than zero).

A blank cell is missing, never zero. \texttt{pandas} reads it as NaN, and the loader's
\texttt{\_without\_missing} turns it into \texttt{None} before any model is built, so a value
that was not measured never arrives as a number that was.

Loading is all or nothing. Every accessor on \texttt{CsvSource} is decorated with
\texttt{@pa.check\_types}, so the table is validated as it is handed out; one offending cell
raises, and no \texttt{Run} object is built at all. That is the right behaviour here because the
alternative is worse: a table loaded with the bad rows quietly dropped would still produce a
polynomial, and it would be the wrong polynomial with no sign that anything had gone missing.
Fixing data is a separate job, which is why \texttt{CsvSource.raw\_tables} exists --- it returns
the tables unchecked, and the \texttt{check} command uses it to list what to repair in the
workbook. Calculating refuses bad data; reporting on bad data does not.

\subsection{The four objects, and how they nest}

\begin{itemize}[nosep]
  \item \texttt{Run} --- the setup, plus two lists: \texttt{stack\_readings} and
        \texttt{tank\_readings}.
  \item \texttt{StackReading} --- the power log at one minute of a run.
  \item \texttt{TankReading} --- one tank at one minute of a run, with an optional
        \texttt{sample}.
  \item \texttt{Sample} --- the lab's results for liquid drawn from one tank at one minute.
\end{itemize}

All four are \emph{pydantic} models, so the nesting is also the validation: handing
\texttt{Run.model\_validate} a dictionary of rows builds and checks the whole tree in one call.

A sample hangs off the tank reading it was drawn from, not off the run. The reason is that a
sample is not a property of a run; it is a property of a moment in a tank. Two tanks are often
sampled at the same minute, so a flat list on the run would need the tank and the minute carried
along to say which bottle belonged where --- and that pairing would then exist twice, once on
the sample and once in whatever lookup used it. Hung off the reading, the pairing is the
structure. \texttt{DataLoader.\_with\_samples} does the matching once, on
\texttt{(tank, minute)}, and after that any code that has a reading has its sample and its
conductivity, pH, temperature and weight together. The end-sample step depends on exactly this:
it walks one tank's readings, keeps the ones that have a sample, sorts them by minute, and looks
at the first and the last.

On the current data, 21 runs carry 1049 tank readings and 891 stack readings, and all 345
samples attach to a reading.

\subsection{The id on every object}

Every object carries an \texttt{id}, and the ids are built so that one string says where a
value came from.

\begin{table}[htbp]
  \centering
  \caption{What each object's \texttt{id} is, with an example from the current data.}
  \label{tab:ids}
  \begin{tabular}{@{}lll@{}}
    \toprule
    Object & \texttt{id} & Example \\
    \midrule
    \texttt{Run}          & the run code                          & \texttt{AJ} \\
    \texttt{TankReading}  & run, tank, minute                     & \texttt{AJ:conc:0} \\
    \texttt{StackReading} & run, minute                           & \texttt{AJ:0} \\
    \texttt{Sample}       & the lab's own label                   & \texttt{AE1} \\
    \texttt{Result}       & the run code                          & \texttt{AE} \\
    \bottomrule
  \end{tabular}
\end{table}

A \texttt{Sample} keeps the lab's label as its id, because that is the string written on the
bottle and the one to quote when querying a result; it carries \texttt{run\_id} and
\texttt{tank\_id} as well, so it can still point back at the reading it came from.
\texttt{TankReading} and \texttt{StackReading} each carry \texttt{run\_id} for the same reason.
\texttt{Result} has an id at all only so that the journal can accept one.

That last point is what the ids are for. \texttt{Journal.\_write} does not care what it has been
handed; it reads \texttt{record.id} and \texttt{type(record).\_\_name\_\_} and writes those down.
So any object with an id can be journalled, and the entry locates the thing without the journal
knowing anything about it. The report then splits the id on colons to sort the entries by run
and then by minute, which is why a dropped run (no minute in its id) prints above its own
readings.

\subsection{The stages and the classes that do them}

\begin{table}[htbp]
  \centering
  \small
  \caption{Each stage, the class that drives it, and what it hands on.}
  \label{tab:stages}
  \begin{tabular}{@{}lll@{}}
    \toprule
    Stage & Class & Hands on \\
    \midrule
    Load       & \texttt{DataLoader}, \texttt{CsvSource}      & \texttt{list[Run]} \\
    Preprocess & \texttt{Preprocessor} and three \texttt{Step}s & \texttt{list[Run]}, shorter \\
    Calculate  & \texttt{Calculator} and the quantity classes & \texttt{list[Result]} \\
    Model      & \texttt{Surrogate}, \texttt{FittedOutcome}, \texttt{DesignPoints}
               & \texttt{SurrogateModels} \\
    \midrule
    Order      & \texttt{Analyzer}  & \texttt{Analysis} \\
    Record     & \texttt{Journal}   & \texttt{tuple[Entry, ...]} \\
    Report     & \texttt{Analyze}   & the printed lines \\
    \bottomrule
  \end{tabular}
\end{table}

\subsection{What follows}

The rest of this document opens each stage in turn: the three preprocessing steps and the six
reading checks, with the thresholds each one judges by and the reason behind every number; the
quantities the calculator works out from tank weights, lab results and the power log, and the
assumptions inside each one; the second-order polynomial, how the two settings are logged and
coded, how terms are kept or dropped, and how leave-one-out error is weighed against the bench's
own scatter; and the search for the best setup under the recovery floor and the voltage ceiling.
